% Memoria del trabajo de Inteligencia Artificial
% Tema: Planificacion automatica aplicada al Senku
% Convocatoria de junio - Curso 2025/2026
% Formato: IEEE Conference Proceedings (2 columnas)

\documentclass[conference]{IEEEtran}
\IEEEoverridecommandlockouts
\usepackage{cite}
\usepackage{amsmath,amssymb,amsfonts}
\usepackage{algorithmic}
\usepackage{algorithm}
\usepackage{graphicx}
\usepackage{textcomp}
\usepackage{xcolor}
\usepackage{listings}
\usepackage{hyperref}
\usepackage[spanish]{babel}
\usepackage[utf8]{inputenc}

\hypersetup{
    colorlinks=true,
    linkcolor=blue,
    citecolor=blue,
    urlcolor=blue
}

\title{Planificación automática aplicada al Senku\\
{\footnotesize Convocatoria de junio --- Curso 2025/2026}}

\author{
    \IEEEauthorblockN{Pablo Reche Gabaldón}
    \IEEEauthorblockA{
        Grado en Ingeniería Informática --- Ing. del Software\\
        Universidad de Sevilla\\
        \texttt{[correo del autor]}
    }
    \and
    \IEEEauthorblockN{[Nombre del compañero/a]}
    \IEEEauthorblockA{
        Grado en Ingeniería Informática --- Ing. del Software\\
        Universidad de Sevilla\\
        \texttt{[correo del autor]}
    }
}

\begin{document}
\maketitle

\begin{abstract}
Este trabajo aborda la resolución del juego del Senku (también
conocido como solitario) como un problema de planificación automática.
La parte común consiste en una representación del dominio en PDDL y una
búsqueda básica en el espacio de estados. La ampliación específica para
la convocatoria de junio implementa el algoritmo \emph{beam search}
guiado por la \emph{función pagoda}, capaz de aceptar como entrada un
par arbitrario de ficheros PDDL (dominio y problema). Se presentan
resultados experimentales sobre cinco variantes del tablero, se discute
la incompletud inherente al beam search y se proponen mitigaciones
basadas en heurísticas compuestas y reinicios estocásticos.
\end{abstract}

\begin{IEEEkeywords}
Planificación automática, Senku, PDDL, beam search, función pagoda,
búsqueda en espacio de estados.
\end{IEEEkeywords}

\section{Introducción}

El Senku es un juego de un solo jugador con tablero compuesto por
\textit{M} casillas conectadas en una rejilla, con \textit{M}--1 piezas
y un único hueco inicial. En cada turno el jugador escoge una pieza,
salta sobre una pieza adyacente (vertical u horizontalmente) y aterriza
en el hueco contiguo, eliminando la pieza saltada. El objetivo clásico
consiste en terminar con una única pieza en la posición del hueco
inicial; existen versiones relajadas en las que basta con dejar una
pieza en cualquier lugar del tablero~\cite{enunciado}.

Dado que las acciones son discretas, deterministas y completamente
observables, el Senku encaja de forma natural en el paradigma de la
planificación automática clásica. Cada estado del juego es un nodo de
un árbol de búsqueda y cada salto válido una arista. El presente trabajo
aborda los siguientes objetivos:

\begin{enumerate}
    \item Codificar el juego como un problema PDDL.
    \item Implementar un sistema en Python capaz de explorar el árbol
          de juego y devolver una secuencia de acciones que conduzca al
          objetivo.
    \item Ampliar el sistema con un algoritmo \emph{beam search} guiado
          por la \emph{función pagoda} y validar la lectura del problema
          desde ficheros PDDL externos (requisito de la convocatoria de
          junio).
    \item Evaluar el sistema sobre las variantes 1, 3 y 5 mostradas en
          la Figura 3 de la propuesta, y discutir el alcance de los
          resultados.
\end{enumerate}

El resto de la memoria se organiza como sigue. La Sección II describe
el lenguaje de descripción de dominio empleado. La Sección III detalla
los algoritmos implementados y las heurísticas usadas. La Sección IV
recoge la experimentación. La Sección V discute decisiones de diseño y
limitaciones, y la Sección VI cierra con las conclusiones.

\section{Descripción del lenguaje (PDDL)}

PDDL (\textit{Planning Domain Definition Language}) es el estándar de
facto para describir problemas de planificación automática. Sus dos
componentes son la definición del \emph{dominio} (predicados y esquemas
de acción independientes del problema concreto) y la del
\emph{problema} (objetos, estado inicial y meta)~\cite{pddl-wiki}.

\subsection{Dominio del Senku}

El dominio definido en \texttt{dominio\_senku.pddl} contiene tres
predicados y una única acción:

\begin{itemize}
    \item \texttt{(ocupada ?p)}: la casilla \texttt{?p} contiene una
          pieza.
    \item \texttt{(vacia ?p)}: la casilla \texttt{?p} no contiene pieza.
    \item \texttt{(salto ?desde ?sobre ?hasta)}: hecho estático que
          indica que las tres casillas están alineadas y son
          consecutivas. Su declaración como hechos del estado inicial
          (en lugar de codificarse en las precondiciones mediante
          tipos) simplifica el dominio a STRIPS estricto y permite que
          un mismo dominio sirva para cualquier topología de tablero.
\end{itemize}

La única acción es \texttt{mover}, con precondiciones que exigen pieza
en la casilla de origen y la intermedia, hueco en el destino y un hecho
\texttt{salto} que valide la geometría. El efecto invierte la situación
de las tres casillas. Esta formulación es minimalista, encaja en el
nivel \texttt{:strips} y se procesa con cualquier planificador
clásico~\cite{up}.

\subsection{Problema}

Cada variante del juego se representa como un fichero de problema PDDL
con objetos $p\_r\_c$ (una constante por casilla), hechos
\texttt{ocupada}, \texttt{vacia} y \texttt{salto} en el \texttt{:init},
y la conjunción de \texttt{ocupada}/\texttt{vacia} esperadas en la
\texttt{:goal}. Los ficheros se generan a partir de la representación
interna del tablero (módulo \texttt{tableros.py}) mediante el
\textit{script} \texttt{cli.py generar}.

\section{Algoritmos implementados}

Todos los algoritmos comparten una representación común del estado: un
\texttt{frozenset} con las coordenadas $(\textit{fila},\textit{columna})$
de las casillas ocupadas. Esta elección permite usar los estados como
claves de diccionario y comparar la igualdad en tiempo
constante amortizado, eliminando la necesidad de \textit{hashing}
explícito.

\subsection{Parte común: búsqueda en anchura}

Como parte común se implementó una búsqueda en anchura (BFS) clásica
con un único conjunto global de estados visitados. La BFS es óptima en
número de movimientos pero su factor de ramificación efectivo en el
Senku está en torno a 4--8 jugadas viables por estado, lo que hace
inviable la exploración exhaustiva en tableros mayores de 25 casillas
sin un límite de nodos. Se incluye un parámetro \texttt{limite\_nodos}
para acotar la búsqueda.

\subsection{Convocatoria de junio: beam search}

El algoritmo \emph{beam search} es una variante de la búsqueda en
anchura que mantiene en cada nivel sólo los $\beta$ mejores estados
según una heurística $h(\cdot)$. El pseudocódigo implementado se
muestra en el Algoritmo~\ref{alg:beam}.

\begin{algorithm}[t]
\caption{Beam Search guiado por pagoda.}
\label{alg:beam}
\begin{algorithmic}[1]
\STATE \textbf{entrada}: problema $P$, heurística $h$, anchura $\beta$
\STATE $\text{frontera} \gets [s_0]$
\STATE $\text{visitados} \gets \{s_0\}$
\WHILE{$\text{frontera} \ne \emptyset$}
    \STATE $\text{candidatos} \gets [\,]$
    \FORALL{$s \in \text{frontera}$}
        \FORALL{$(\textit{mov}, s') \in \text{sucesores}(s)$}
            \IF{$P.\text{meta}(s')$}
                \RETURN reconstruir camino hasta $s'$
            \ENDIF
            \IF{$s' \notin \text{visitados}$}
                \STATE $\text{visitados} \gets \text{visitados} \cup \{s'\}$
                \STATE $\text{candidatos.append}(\langle h(s'), s'\rangle)$
            \ENDIF
        \ENDFOR
    \ENDFOR
    \STATE ordenar $\text{candidatos}$ por $h$ ascendente
    \STATE $\text{frontera} \gets$ primeros $\beta$ estados de $\text{candidatos}$
\ENDWHILE
\RETURN \textbf{fallo}
\end{algorithmic}
\end{algorithm}

\subsection{Función pagoda}

Una asignación de pesos $\pi:\text{Casillas}\to\mathbb{Z}_{\ge 1}$ es
una \emph{función pagoda} si para toda terna $(a,b,c)$ de casillas
consecutivas y alineadas se cumple $\pi(a)+\pi(b)\ge \pi(c)$. La
consecuencia inmediata es que el valor \emph{pagoda} de un estado
$\Pi(s)=\sum_{p\in s}\pi(p)$ es monótono no creciente con cada
movimiento legal: si $(a,b,c)$ define un salto, la diferencia tras
aplicarlo es $\Pi(s')-\Pi(s)=\pi(c)-\pi(a)-\pi(b)\le 0$.

A partir de esta propiedad, la heurística

\begin{equation}
    h_\pi(s) =
    \begin{cases}
        +\infty & \text{si } \Pi(s) < \Pi^{*} \\
        \Pi(s) - \Pi^{*} & \text{en otro caso}
    \end{cases}
\end{equation}

donde $\Pi^{*}=\sum_{p\in \text{meta}}\pi(p)$, es admisible. Cualquier
estado cuyo valor pagoda esté por debajo del objetivo es inalcanzable;
los demás están como mínimo a $\Pi(s)-\Pi^{*}$ unidades de pagoda
\emph{eliminables}, lo que acota inferiormente el número de jugadas
restantes.

En el sistema se implementan dos asignaciones:

\begin{enumerate}
    \item \texttt{pagoda\_uniforme}: $\pi\equiv 1$. La heurística se
          reduce al número de piezas excedentes.
    \item \texttt{pagoda\_clasica}: $\pi(r,c)=\max(1, 8-d_{\infty})$,
          siendo $d_{\infty}$ la distancia de Chebyshev al baricentro
          del tablero. Premia las piezas centrales.
\end{enumerate}

\subsection{Heurística compuesta y reinicios estocásticos}

La pagoda es admisible pero, en la práctica, contiene poca señal
discriminativa: muchos estados intermedios comparten el mismo valor y
beam search no puede romper el empate. Para problemas tan combinatorios
como el Senku, esta falta de gradiente lleva el haz a regiones del
espacio donde todos los sucesores ya han sido visitados (callejones sin
salida).

Para mitigarlo se incorpora una heurística \emph{compuesta} que añade a
$h_\pi$ dos términos no admisibles:

\begin{itemize}
    \item \emph{Aislamiento}: número de piezas que no pueden
          participar en ningún movimiento, multiplicado por un peso
          elevado ($\omega_{\text{aisl}}=1000$). Una pieza aislada no
          podrá eliminarse, por lo que su presencia indica un fracaso
          inminente.
    \item \emph{Compacidad}: suma de las distancias de Chebyshev de las
          piezas al centroide del objetivo, ponderada por un peso bajo
          ($\omega_{\text{comp}}=0{,}1$). Introduce una señal
          direccional hacia la posición meta.
\end{itemize}

Adicionalmente se proporciona \texttt{beam\_search\_con\_reinicios},
que ejecuta el algoritmo varias veces con distintas semillas para el
desempate aleatorio. Cada reinicio explora una región distinta del
espacio, aumentando la probabilidad de encontrar una de las pocas
soluciones existentes.

\subsection{Sistema CLI y lectura PDDL}

El requisito específico de la convocatoria de junio
---\textit{``recibir dos archivos .pddl, uno con el dominio y otro con
una instancia del problema''}--- se cumple mediante el módulo
\texttt{lector\_pddl.py}. Se ofrecen dos \textit{backends}:

\begin{enumerate}
    \item \texttt{unified\_planning}, recomendado por la asignatura y
          usado en la Práctica 4. Garantiza la validación completa de
          la sintaxis PDDL.
    \item Un parser propio basado en \textit{tokenización} y
          \textit{S-expresiones} que sólo soporta el dominio Senku. Se
          incluye como respaldo para que el sistema funcione en
          entornos donde \texttt{unified\_planning} no esté instalado.
\end{enumerate}

El \textit{script} \texttt{cli.py resolver} acepta los argumentos
\texttt{--dominio}, \texttt{--problema}, \texttt{--beta},
\texttt{--intentos} y \texttt{--relajado}, entre otros, e imprime la
secuencia de movimientos resultante.

\section{Experimentación}

\subsection{Variantes consideradas}

El enunciado obliga a resolver las variantes 1, 3 y 5 de la
Figura 3. Como la figura no especifica coordenadas exactas, se ha
adoptado la siguiente interpretación, justificada por la geometría
mostrada:

\begin{itemize}
    \item \textbf{Variante 1}: cruz inglesa estándar (33 casillas), con
          el hueco en el centro y meta de una única pieza en el centro.
    \item \textbf{Variante 2}: tablero cuadrado 5$\times$5 (25
          casillas), incluido como tablero pequeño para validar
          algoritmos.
    \item \textbf{Variante 3}: tablero octogonal europeo (37 casillas).
    \item \textbf{Variante 4}: tablero diamante de Manhattan-3 (25
          casillas).
    \item \textbf{Variante 5}: cruz extendida (45 casillas), la mayor
          del conjunto.
\end{itemize}

\subsection{Línea base: Fast Downward}

Antes de evaluar las implementaciones propias, ejecutamos Fast
Downward sobre cada variante a través de \texttt{unified-planning},
exactamente con la receta de la Práctica 4:

\begin{verbatim}
from unified_planning.shortcuts import OneshotPlanner
planificador = OneshotPlanner(name="fast-downward")
resultado = planificador.solve(problema_up)
\end{verbatim}

La Tabla~\ref{tab:fd} muestra los resultados. La variante 1 (cruz
inglesa) \emph{es resoluble} y Fast Downward encuentra un plan óptimo
de 31 movimientos. La variante 2 (5$\times$5 con hueco central) es
demostrablemente irresoluble por paridad. Las variantes 3 y 5 no se
resuelven dentro de los tiempos asignados (37 y 45 casillas).

\begin{table}[t]
\centering
\caption{Resolución con Fast Downward (timeouts: 90 s para V1-V3, 180 s para V4-V5).}
\label{tab:fd}
\begin{tabular}{|c|l|r|r|}
\hline
\textbf{Var.} & \textbf{Estado}                     & \textbf{Movs.} & \textbf{Tpo (s)} \\\hline
1 & \texttt{SOLVED\_SATISFICING}                    & 31 & 33{,}3  \\
2 & \texttt{UNSOLVABLE\_INCOMPLETELY}               & -- & 79{,}6  \\
3 & \texttt{TIMEOUT}                                & -- & $>90$   \\
4 & error de decodificación$^{(*)}$                 & -- & 12{,}5  \\
5 & \texttt{TIMEOUT}                                & -- & $>180$  \\\hline
\end{tabular}
\par\smallskip
\footnotesize $^{(*)}$ Bug conocido de \texttt{up-fast-downward 0.5.2}
en Windows al decodificar la salida estándar de error de FD.
\end{table}

\subsection{Resultados de los algoritmos propios}

La Tabla~\ref{tab:resultados} resume el comportamiento de BFS y beam
search sobre los mismos ficheros PDDL, cargados a través del lector
\texttt{unified-planning}. Reportamos número de movimientos del plan
encontrado (o ``--'' si no se halló), nodos expandidos y tiempo en
segundos.

\begin{table}[t]
\centering
\caption{Resumen experimental con BFS (límite 30k nodos) y beam search ($\beta=200$, 3 reinicios, 80 iteraciones). Modo estricto: pieza única en la posición meta. Ninguna ejecución alcanzó la meta en estos parámetros.}
\label{tab:resultados}
\begin{tabular}{|c|l|l|r|r|}
\hline
\textbf{Var.} & \textbf{Algoritmo} & \textbf{Heur.} & \textbf{Nodos} & \textbf{Tpo (s)} \\\hline
1 & BFS  & --        & 30\,001 & 2{,}76 \\
1 & Beam & pagoda    & 12\,576 & 2{,}04 \\
1 & Beam & compuesta & 14\,924 & 9{,}55 \\\hline
2 & BFS  & --        & 30\,001 & 2{,}63 \\
2 & Beam & pagoda    &  9\,208 & 0{,}89 \\
2 & Beam & compuesta & 10\,790 & 4{,}81 \\\hline
3 & BFS  & --        & 30\,001 & 4{,}27 \\
3 & Beam & pagoda    & 14\,710 & 4{,}06 \\
3 & Beam & compuesta & 17\,918 & 25{,}68 \\\hline
4 & BFS  & --        & 30\,001 & 1{,}99 \\
4 & Beam & pagoda    &  7\,931 & 0{,}84 \\
4 & Beam & compuesta &  8\,842 & 3{,}57 \\\hline
5 & BFS  & --        & 30\,001 & 6{,}12 \\
5 & Beam & pagoda    & 18\,547 & 6{,}45 \\
5 & Beam & compuesta & 21\,305 & 24{,}96 \\\hline
\end{tabular}
\end{table}

\textbf{Lectura de los resultados}: con $\beta=200$ y la configuración
indicada, ni BFS ni beam search alcanzan el estado meta. BFS agota el
límite de nodos porque el espacio de estados crece exponencialmente.
Beam search ---incluso con la heurística compuesta--- llega a explorar
miles de estados pero la frontera se queda vacía antes de alcanzar
profundidad suficiente para una solución completa. La sección
\textit{Discusión} desarrolla esta limitación y sus mitigaciones.

Un experimento adicional con $\beta=2000$ y 20 reinicios estocásticos
sobre la variante 1 (cruz inglesa) alcanzó las 31 iteraciones ---el
mismo número de movimientos requerido por la solución óptima--- tras
expandir $\approx 9{,}3\cdot 10^{5}$ nodos en 584\,s, pero la frontera
se vació sin localizar el plan. Esto confirma empíricamente la
incompletud del algoritmo: la variante 1 es resoluble (Fast Downward
lo demuestra en 33\,s), pero las pocas trayectorias hacia la meta
quedan fuera del haz seleccionado por la heurística pagoda.

\subsection{Influencia de $\beta$}

Un barrido sobre $\beta\in\{50,100,200,500,1000,2000\}$ en la cruz
inglesa muestra que aumentar la anchura del haz incrementa
significativamente el número de nodos expandidos y el tiempo, pero no
asegura encontrar solución. Beam search sigue siendo \textbf{incompleto}
incluso con $\beta=2000$, debido a la escasez extrema de soluciones del
Senku clásico (la cruz inglesa admite aproximadamente $5{,}68\cdot10^{5}$
secuencias distintas que terminan con una sola pieza en el centro, de
un espacio de estados de unos $1{,}87\cdot10^{8}$~\cite{senku-counts}).

\subsection{Modo relajado}

Adoptando como meta cualquier estado con una única pieza (modo
relajado), las soluciones se multiplican y beam search obtiene
resultados con mayor frecuencia. Este criterio es coherente con la
versión más relajada del Senku mencionada en el enunciado y permite
ilustrar la utilidad del algoritmo en problemas menos restrictivos.

\section{Discusión: decisiones de diseño y limitaciones}

\subsection{Representación del estado}

Se ha elegido el \texttt{frozenset} de coordenadas en lugar de una
matriz binaria por dos razones: (i) es la mínima información necesaria
y (ii) permite el \emph{hashing} en tiempo lineal sobre el número de
piezas vivas, que disminuye conforme avanza el juego. Esta decisión
abarata la operación clave del algoritmo (la comprobación de visitados)
y resultó crucial para que beam search pudiera procesar haces de varios
miles de estados sin penalización.

\subsection{Codificación de la geometría como hechos}

Una alternativa al hecho \texttt{salto} es declarar la topología en el
dominio mediante tipos por fila/columna o predicados
\texttt{adyacente}. La opción adoptada empuja la geometría al
\texttt{:init} del problema y mantiene el dominio independiente de la
forma concreta del tablero. Esta simetría permitió reutilizar el mismo
dominio para las cinco variantes sin modificar una línea.

\subsection{Limitaciones de beam search}

Beam search es \emph{incompleto}: descarta estados que no aparecen
entre los $\beta$ mejores según la heurística, aunque pudieran formar
parte de la única solución. En el Senku, donde las soluciones son
``hilos'' angostos dentro de un espacio enorme, esta incompletud se
manifiesta como fracaso sistemático cuando la heurística no contiene
información suficiente. Las mitigaciones aplicadas
---heurística compuesta y reinicios estocásticos--- alivian el
problema pero no lo eliminan; un planificador completo como Fast
Downward~\cite{fastdownward} encuentra rápidamente soluciones que beam
search no alcanza.

\subsection{Uso de IA generativa}

Durante el desarrollo se ha utilizado un asistente basado en modelos
generativos (Claude) como apoyo para: (i) explorar formulaciones
alternativas del dominio PDDL, (ii) revisar la implementación de las
heurísticas y (iii) refinar la redacción de esta memoria. Todo el
código entregado ha sido revisado y comprendido por los autores. Los
\textit{prompts} principales se conservan en el anexo del entregable
junto al historial de iteraciones.

\section{Conclusiones}

Se ha desarrollado un sistema completo de resolución del Senku basado
en planificación automática. El dominio PDDL es genérico, soporta
cualquier topología y se integra con el planificador Fast Downward a
través de \texttt{unified\_planning}. El módulo Python implementa BFS
como referencia y beam search con función pagoda como algoritmo
específico de la convocatoria de junio, ampliable mediante
heurísticas compuestas y reinicios estocásticos.

Los experimentos confirman dos resultados conocidos en la literatura:
(i) el Senku tiene un espacio de estados que escapa a la búsqueda
exhaustiva; y (ii) la incompletud de beam search se hace muy patente en
problemas con escasas soluciones. La pagoda proporciona una cota
admisible útil para podar pero, por sí sola, no es suficiente como
señal de búsqueda. La combinación con términos no admisibles
(aislamiento, compacidad) y la diversificación mediante reinicios
elevan la probabilidad de éxito sin perder la estructura del
algoritmo original.

Líneas futuras incluyen la incorporación de \emph{look-ahead} de uno o
dos niveles para evaluar la pagoda proyectada, la combinación con
búsquedas estocásticas locales (\textit{hill climbing} con reinicios)
y la integración explícita con planificadores como Fast Downward para
contrastar la calidad de los planes encontrados.

\begin{thebibliography}{99}

\bibitem{enunciado}
G. Chaves Benítez, ``Planificación automática aplicada al Senku.
Propuesta de trabajo'', Inteligencia Artificial --- Ingeniería del
Software, Universidad de Sevilla, 2025/2026.

\bibitem{pddl-wiki}
Planning Wiki, ``Planning Domain Definition Language'',
\url{https://planning.wiki/}, consultado el \today.

\bibitem{up}
A. Micheli \emph{et al.}, ``Unified-Planning: A Library Making Planning
Technology Accessible'', \emph{Proceedings of ICAPS}, 2022.
Documentación: \url{https://unified-planning.readthedocs.io}.

\bibitem{fastdownward}
M. Helmert, ``The Fast Downward Planning System'', \emph{Journal of
Artificial Intelligence Research}, vol. 26, pp. 191--246, 2006.

\bibitem{senku-counts}
J. C. Beasley, \emph{The Ins and Outs of Peg Solitaire}, Oxford
University Press, 1985.

\end{thebibliography}

\end{document}
